\chapter{A Demonstrator for Self-organizing Robot Teams}
\label{chap:demonstrator}

% NOTE: this chapter reports Project Emerge as the physical realisation of one
% deployment admitted by the model of \cref{chap:pulverization}. The pulverisation
% reading of its architecture is the thesis's own: the paper never uses that
% vocabulary. Aggregate computing background lives in
% \cref{subsec:aggregate-computing}; reference it instead of re-explaining.
% Two caveats must stay visible: the runtime exchanges neighbour messages in
% memory rather than over forwarding chains (\cref{sec:emerge-as-pulverisation}),
% and the chapter does not exercise
% \crefrange{chap:language-and-coordination}{chap:heterogeneous-gnn} --
% \cref{sec:emerge-discussion} states that gap rather than claiming otherwise.

Every claim this thesis has made about deployment has been checked in simulation:
the pulverisation model in Alchemist (\cref{subsec:pulverization-validation}),
the reconfiguration policies of \cref{chap:deployment-and-reconfiguration} in the same simulator,
and the learned offloading policy of \cref{chap:heterogeneous-gnn} against a simulated continuum.
%
Simulation is the right instrument for the questions those chapters ask, because it makes topologies, failures, and network sizes controllable.
%
It leaves one question open: whether a deployment the model admits can be built on real devices, with the sensing and actuation those devices have.

This chapter reports a system that answers it.
%
\emph{Project Emerge} is a toolchain that runs a single aggregate program on a team of physical mobile robots.
It was demonstrated publicly at the European Researchers' Night in 2024 and 2025.
%
Nine robots form spatial patterns and recover them after one of them is displaced by hand.
%
This chapter is concerned with where the computation runs.
The robots carry an ESP32 microcontroller and no aggregate runtime,
while a server evaluates the program on their behalf and returns actuation commands.
%
In the terms of \cref{chap:pulverization}, that arrangement is a deployment, with the robots as owners and the server as their shared surrogate.

That correspondence is the argument of this chapter. The paper does not make it.
%
Project Emerge was built to make aggregate computing demonstrable on hardware, and it describes its own architecture as centralised, without reference to pulverisation.
%
The deployment it arrived at, for reasons of its own -- microcontrollers too small to host a runtime -- is one of the deployments licensed by \cref{thm:deployment-independence},
and the freedom the theorem grants would let the same program run elsewhere unchanged.
%
The model left the deployment ``taken as given'' (\cref{sec:pulverization-discussion}); this chapter builds one such deployment on hardware.

The chapter follows the demonstrator from requirement to discussion:
what running a macro-program on real robots demands of the substrate (\cref{sec:emerge-requirements}),
the toolchain built to meet those demands (\cref{sec:emerge-architecture}),
how its deployment reads in the vocabulary of \cref{chap:pulverization} (\cref{sec:emerge-as-pulverisation}),
what the robots did (\cref{sec:emerge-demonstration}),
and what the demonstrator establishes and leaves untouched (\cref{sec:emerge-discussion}).
%
The contribution is published as~\cite{Farabegoli:2026:ProjectEmerge}, extending~\cite{DBLP:conf/coordination/AguzziBBCCDFPV25}.

\section{What a Physical Deployment Requires}
\label{sec:emerge-requirements}

Aggregate computing assumes a device that senses, computes, and actuates in rounds (\cref{subsubsec:ac-execution-model}).
%
A simulator supplies all three by construction.
On physical robots each has to be built, or substituted by something that stands in for it,
and a demonstrator's engineering effort goes into closing that gap.
%
Six requirements describe the gap.

\paragraph{Positioning.}
A macro-program that computes movement needs to know where its devices are.
%
Position alone is not enough:
a command to move in a direction is meaningless to a robot that does not know which way it is facing,
so the positioning system must report \emph{orientation} as well.
%
Outdoors this can be satellite positioning; indoors, where the demonstrator runs, it cannot.

\paragraph{Homogeneous program loading.}
One program specifies the behaviour of the whole team (\cref{sec:macro-programming}),
and every robot must be running it for the collective behaviour to be the one specified.
%
Loading it must not require shutting the team down.
%
During a switch from one version to another the two coexist,
and they must not interfere while they do, until every robot has moved across.

\paragraph{Neighbourhood policy.}
A device's neighbours are a dynamic set, and in a physical deployment membership is decided by physical facts:
proximity, line of sight, achievable data rate.
%
The substrate has to make that policy explicit and replaceable, rather than fixing one notion of adjacency.

\paragraph{Actuation agnosticism.}
The output of a round must not presuppose a chassis.
%
The program's output is an intention -- move this way, turn that way --
and translating an intention into wheel or rotor commands belongs to the platform, not to the program.
%
Without that separation the same program cannot run on two different robots.

\paragraph{Openness to platforms, and to where the program runs.}
The previous requirement concerns which robots a program can drive;
this one concerns which robots can drive themselves.
%
Platforms differ in computational capacity,
so the substrate must be able to evaluate the program on the robot, or remotely on a supporting host that does it on the robot's behalf.
%
This requirement matters most for this thesis,
and the demonstrator states it but satisfies only its remote half (\cref{sec:emerge-as-pulverisation}).

\paragraph{Interaction.}
Emergent behaviour is hard to trust when there is nothing to look at.
%
An operator needs to observe device states and the fields the program computes,
and to perturb the running system -- displace a device, change a parameter, swap the program -- while it runs.

\section{Architecture of the Demonstrator}
\label{sec:emerge-architecture}

Project Emerge meets those six requirements with five subsystems communicating over a single MQTT broker:
a localisation system, a neighbourhood manager, the aggregate runtime, a dashboard, and the robots themselves.
%
\Cref{fig:emerge-architecture} shows how they connect.
%
Because every exchange goes through one broker, a simulator can later replace the robots (\cref{sec:emerge-demonstration}):
nothing in the runtime knows whether the positions it reads come from a camera or from a simulation.

\begin{figure}[t]
    \centering
    \includegraphics[width=0.85\textwidth]{emerge-architecture.pdf}
    \caption[System-level architecture of the demonstrator]{System-level architecture of the demonstrator. The localisation system tracks robot positions, the neighbourhood system maintains the adjacency relation, the aggregate runtime evaluates the collective program, and the dashboard supports monitoring and perturbation. Solid arrows are message flows through the broker; dashed arrows are logical dependencies. Adapted from~\cite{Farabegoli:2026:ProjectEmerge}.}
    \label{fig:emerge-architecture}
\end{figure}

Localisation is a camera above the arena and an ArUco marker on each robot,
processed by an OpenCV-based detector that publishes each robot's identifier, position, and orientation.
%
This discharges the positioning requirement without putting any sensor on the robot, a substitution examined in \cref{sec:emerge-as-pulverisation}.
%
The neighbourhood manager derives adjacency from those positions by spatial proximity,
and is kept separate from the localisation component so that a different policy can replace it.
%
The dashboard renders robots, their neighbourhood relation, and the current program,
and lets an operator drag a robot to a new position, retune a parameter, elect a leader, or load a different program.

\subsection{The Aggregate Runtime}
\label{subsec:emerge-runtime}

The runtime is written in Scala on ScaFi (\cref{subsec:cas-platforms}),
with the collective behaviours taken from MacroSwarm,
a library of field-based building blocks for swarm behaviour -- flocking, leader-based coordination, shape formation -- built on top of ScaFi~\cite{DBLP:conf/coordination/AguzziCV23}.
%
It is factored into four abstractions plus an entry point, and that factoring carries the requirements of \cref{sec:emerge-requirements} into code.\footnote{The toolchain, the companion rover hardware, and a reproducible capsule are released as open source: \url{https://github.com/Project-Emerge} and \url{https://doi.org/10.5281/zenodo.18224119}.}

An \texttt{Environment} models the physical context of one evaluation:
which robots are currently present, where each is, who its neighbours are, and what the sensors read.
%
An \texttt{EnvironmentProvider} builds that snapshot on demand,
its broker-based implementation assembling it from the localisation and neighbourhood messages received so far.
%
An \texttt{EnvironmentUpdater} performs the reverse translation, taking what a round computed and delivering it to the robots.
%
An \texttt{Orchestrator} evaluates the program:
given an environment, it computes, \emph{for each robot in it}, the program that robot is to run, evaluates it, and returns that robot's command.
%
The \texttt{AggregateService} composes the three into a loop that samples an environment, evaluates, and dispatches, at a fixed frequency.
%
\Cref{fig:emerge-execution-cycle} shows one turn of that loop.

\begin{figure}[t]
    \centering
    \includegraphics[width=0.6\textwidth]{emerge-execution-cycle.pdf}
    \caption[One evaluation cycle of the aggregate runtime]{One evaluation cycle of the aggregate runtime. The environment provider builds a logical network model from the physical robot states at time $t$; the orchestrator evaluates the aggregate program over it, producing actuation commands; the updater dispatches them; the robots execute them, carrying the system to $t+1$. Adapted from~\cite{Farabegoli:2026:ProjectEmerge}.}
    \label{fig:emerge-execution-cycle}
\end{figure}

Because the loop runs at a fixed frequency, it needs no explicit treatment of message ordering or timestamps:
each cycle consumes the latest message available on each topic and discards older ones.
%
A robot whose position update has not arrived is therefore evaluated against a stale pose rather than none.
%
This is a source of error the model already accounts for, since \cref{thm:deployment-independence} treats delayed and lost messages as transient perturbations that self-stabilising components absorb.

The boundary between program and platform is a small set of commands, shown in \cref{lst:emerge-actuation}:
rotate by a vector, move forward along a vector, hold, stop.
%
A MacroSwarm program yields one of these rather than a motor setting --
a prescriptive action, in the vocabulary of \cref{subsubsec:ac-execution-model}, that the robot is expected to attempt.

\noindent\begin{minipage}{\linewidth}
\begin{lstlisting}[
  style=scala3,
  caption={The boundary between an aggregate program and the robot platform. A program produces an intention; nothing in the type mentions wheels, rotors, or a chassis.},
  label={lst:emerge-actuation},
]
enum Actuation:
  case Rotation(towards: Vector2D)
  case Forward(along: Vector2D)
  case NoOp
  case Stop

trait BaseDemo extends AggregateProgram:
  override def main(): Actuation
\end{lstlisting}
\end{minipage}

The updater turns an intention into motion,
and in the current system it does so for differential drive:
a desired velocity vector becomes a pair of left and right wheel velocities.
%
Supporting a different chassis means writing a different updater; the program is unchanged.
%
The default program, \texttt{PointTheLeader}, shows the same boundary from the program's side:
it propagates a vector field outward from a leader elected through the dashboard,
using a gradient cast,
and returns a \texttt{Rotation} along that field so each robot turns to face the leader.

The runtime loads one program for the whole team, and can be told to switch to another over the broker while the team is running.
%
Because every subsystem speaks only to the broker,
a simulation engine can stand in for the robots and the camera together, publishing the same messages a physical arena would.

\section{The Deployment as a Pulverised One}
\label{sec:emerge-as-pulverisation}

The description just given is the demonstrator's own.
%
Restated in the terms of \cref{chap:pulverization}, it describes a deployment, subject to one qualification, stated below.
%
Nothing in this section is a claim the demonstrator makes about itself.

The nine robots are the application devices $\deviceIdSet$:
they carry the actuators the macro-program drives, and each of them owns an instance of its behaviour.
%
None of them executes that instance.
%
The server hosting the runtime is an infrastructural device, and it is the surrogate for every robot's behaviour component --
the case anticipated in \cref{subsec:forwarding-chains}, in which many owners forward the same component to one shared host.
%
The model permits it because component instances hold no state beyond their round-inputs,
which is why a surrogate can serve many owners as one without losing owner-specific behaviour.
%
The \texttt{Orchestrator} implements that collapse:
one process, one instance per robot, each evaluated against its own view of the network.
%
The broker relays messages without executing any, the role of the intermediate hop $\deviceId_1$ in \cref{fig:pulverization-forwarding}.
%
Writing the deployment predicate of \cref{subsec:forwarding-chains} for robot $\deviceId_j$ gives
%
\[
    \deploy(\deviceId_j,\ \component_{\mathit{beh}},\ \deviceId_j\ \mathit{broker}\ \mathit{server}),
\]
%
with actuation left at the owner.
%
Every pair of a device and a component is deployed and no chain repeats a device, so the deployment is correct in the model's sense.

One qualification limits what the restatement covers.
%
The broker carries the sensing and actuation traffic, but not the neighbour messages:
the runtime exchanges those in memory, inside the single process that hosts every instance, avoiding a network round trip per neighbour per round.
%
Under the model those messages would travel along forwarding chains.
%
The demonstrator therefore realises the limit case in which every chain terminates at the same host,
so the exchange the model would put on the network happens inside the surrogate instead.
%
The model admits that case -- it is the same collapsing of instances licensed in \cref{subsec:forwarding-chains} --
and the choice behind it is about performance, not about semantics.
%
A deployment that spread the same instances over several surrogates would have to send those messages, and would pay for them.

Sensing is the second place where the model's hypotheses constrain the design.
%
The robots' onboard sensors were not used:
they were replaced by the overhead camera, whose estimates of position and orientation reach the runtime through the environment provider.
%
The choice was made to keep a public demonstration reliable,
and in the model's terms it is a sensing component that does not run at its owner either.
%
A surrogate reads its own sensors, not the owner's,
so the substitution is sound only if what it samples is indistinguishable, as far as the component is concerned, from what the owner would have sampled --
the requirement that sensor state be uniform along the forwarding chain (\cref{subsec:deployment-independence}).
%
The requirement is met by construction:
the camera reports each robot's own position and orientation, the same quantities an onboard sensor would have reported for that robot.
%
The demonstrator thus satisfies the theorem's second hypothesis for the one sensing modality it uses,
and it satisfies the first by building its behaviours out of MacroSwarm's self-stabilising blocks (\cref{subsec:aggregate-computing}).

The arrangement the demonstrator settled on is not an alternative to the pulverisation model but an instance of it, arrived at independently,
and the reasons that forced it -- microcontrollers that cannot host a runtime, sensing better done from the ceiling than from the chassis --
are exactly the capability mismatches for which offloading is the answer given in \cref{subsec:forwarding-chains}.
%
The requirement of \cref{sec:emerge-requirements} that the program be executable \emph{either} on the robot \emph{or} on a supporting host
names the same freedom \cref{thm:deployment-independence} grants.
%
The demonstrator states that requirement and meets its remote half only.
%
Its own plan for the other half -- each robot running a local instance of the runtime, with localisation decentralised through radio ranging --
is, in the model's terms, the monolithic deployment in which every component runs at its owner.
%
The two are the extremes of one space; the mixed deployments lie between them,
and revising the choice among them at run time is the problem addressed in \crefrange{chap:deployment-and-reconfiguration}{chap:heterogeneous-gnn}.
The demonstrator does not yet do so.

\section{Demonstration and Observed Behaviour}
\label{sec:emerge-demonstration}

The demonstrator ran twice as a public exhibit, at the European Researchers' Night in 2024 and 2025,
in an indoor venue with an untrained audience walking through it.
%
The team was nine mobile robots on chassis designed and printed in-house,
each carrying an ESP32 microcontroller responsible for communication and low-level actuation, and a marker for the camera to track.

Five programs were demonstrated, chosen to exercise different coordination requirements:
a circle, in which the robots distribute along a circumference around a leader elected at run time;
a square, in which they converge on four corners while keeping the edges aligned;
a V-shape, two symmetric branches from a leader;
point-the-leader, in which each robot orients towards a designated one;
and line formations, horizontal and vertical, with uniform spacing along an axis.
%
Between them they cover symmetry, alignment, leader-based coordination, and global geometric constraints.
%
\Cref{fig:emerge-formations} shows four of them as they appeared in the venue.

\begin{figure}[tbp]
    \centering
    \begin{subfigure}{0.47\textwidth}
        \centering
        \includegraphics[width=\textwidth]{emerge-formation-circle.jpg}
        \caption{Circle}
    \end{subfigure}
    \hfill
    \begin{subfigure}{0.47\textwidth}
        \centering
        \includegraphics[width=\textwidth]{emerge-formation-square.jpg}
        \caption{Square}
    \end{subfigure}

    \vspace{0.6em}

    \begin{subfigure}{0.47\textwidth}
        \centering
        \includegraphics[width=\textwidth]{emerge-formation-vshape.jpg}
        \caption{V-shape}
    \end{subfigure}
    \hfill
    \begin{subfigure}{0.47\textwidth}
        \centering
        \includegraphics[width=\textwidth]{emerge-formation-pointleader.jpg}
        \caption{Point-the-leader}
    \end{subfigure}
    \caption[Spatial patterns formed by the robot team]{Four of the spatial patterns the team formed during the demonstration, each produced by a different aggregate program loaded into the running system. From~\cite{Farabegoli:2026:ProjectEmerge}.}
    \label{fig:emerge-formations}
\end{figure}

Programs were switched while the team was running, as the requirement of \cref{sec:emerge-requirements} demands.
The audience saw one formation dissolve into the next without the system being restarted.
%
The dashboard, shown in \cref{fig:emerge-dashboard}, was part of the exhibit rather than an operator tool,
since the neighbourhood relation it draws is otherwise invisible,
and a visitor could change the neighbourhood radius and watch the formation reorganise.

\begin{figure}[tbp]
    \centering
    \includegraphics[width=\textwidth]{emerge-dashboard.png}
    \caption[The dashboard used to observe and perturb the running team]{The dashboard used to observe and perturb the running team. Lines between robots draw the current neighbourhood relation, which is otherwise invisible to an observer. The controls load a different aggregate program and retune the neighbourhood radius; selecting a robot exposes its state, a manual movement control, and the option to elect it leader. From~\cite{Farabegoli:2026:ProjectEmerge}.}
    \label{fig:emerge-dashboard}
\end{figure}

The property worth testing on hardware was self-stabilisation.
%
Simulation had already confirmed that an offloaded deployment converges where a monolithic one does (\cref{subsec:pulverization-validation}),
under the perturbations that study modelled: node mobility, and movement stopping partway through.
%
A physical arena admits perturbations neither of those covers, applied by hand and not scripted in advance.
%
Formations were disrupted in three ways:
displacing one or more robots mid-execution,
removing a robot from the arena and later putting it back,
and changing the formation parameters or the elected leader while the team was converging.
%
Removal and reinsertion is the most informative of the three, because it is node churn --
one of the perturbation classes the definition of self-stabilisation quantifies over (\cref{subsec:deployment-independence}) --
applied to hardware rather than to a simulated network.
%
In every case tried, the team absorbed the perturbation and recovered the formation without intervention.
%
\Cref{fig:emerge-selfhealing} shows one such recovery.

\begin{figure}[tbp]
    \centering
    \begin{subfigure}{0.32\textwidth}
        \centering
        \includegraphics[width=\textwidth]{emerge-selfhealing-1.png}
        \caption{Formation reached}
    \end{subfigure}
    \hfill
    \begin{subfigure}{0.32\textwidth}
        \centering
        \includegraphics[width=\textwidth]{emerge-selfhealing-2.png}
        \caption{A robot displaced by hand}
    \end{subfigure}
    \hfill
    \begin{subfigure}{0.32\textwidth}
        \centering
        \includegraphics[width=\textwidth]{emerge-selfhealing-3.png}
        \caption{Formation recovered}
    \end{subfigure}
    \caption[Recovery of a line formation after manual perturbation]{Recovery of a line formation after a robot is displaced by hand during execution. From~\cite{Farabegoli:2026:ProjectEmerge}.}
    \label{fig:emerge-selfhealing}
\end{figure}

The evidence is qualitative.
%
The demonstration measured nothing.
It supports the claim that self-stabilisation survives physical embodiment and manual perturbation,
and no claim about how quickly, or about how recovery time compares across deployments.
%
A quantitative version of this experiment is within reach through the simulation engine that replaces the robots and the camera:
it exposes the same broker interface,
so a program validated against it and a program run in the arena are the same program under two sensing substrates.

\section{Discussion}
\label{sec:emerge-discussion}

The demonstrator establishes something narrow.
%
A deployment admitted by the model of \cref{chap:pulverization} has been built on real hardware and run in front of an audience.
%
That is evidence that the model's admissible deployments are not only simulable but realisable,
and that the capability mismatches motivating offloading in \cref{subsec:forwarding-chains} are the ones physical robots present.
%
The sensing substitution shows this most clearly:
condition \emph{(ii)} of \cref{thm:deployment-independence} reads as a technical hypothesis on paper,
and here it is a design decision -- put the sensor on the ceiling instead of the chassis -- that happens to satisfy it.

The demonstrator leaves much of the thesis untested.
%
Its deployment is fixed by hand before the system starts and never revised while it runs,
so it exercises none of \crefrange{chap:deployment-and-reconfiguration}{chap:heterogeneous-gnn}:
there is no reconfiguration policy, no declarative plan, and no learned offloading decision in it.
%
Its behaviour is one aggregate program rather than a macro-program decomposed into a \ac{DAG} of components,
so the finer-grained partitioning of \cref{subsec:macroprogram-dag} is not put to work either;
what is offloaded is the whole behaviour, as in the original five-role pulverisation.
%
It is written in ScaFi and MacroSwarm directly,
not through the unified capabilities of \cref{chap:language-and-coordination} nor with the generation support of \cref{chap:language-based-macroprogramming}.
%
And it reports no measurements.
%
The demonstrator is therefore a realisation of one point of the deployment space this thesis defines, on physical devices,
rather than end-to-end validation of the thesis as a whole.

Those gaps also make it the right place to exercise the rest.
%
The requirement it states and half-satisfies -- run the program on the robot or on a supporting host --
is a deployment choice, and its architecture already isolates that choice behind the environment provider and the updater.
The platform this thesis proposes would let a policy make that choice instead of an engineer.
%
Four steps would close the distance, each one connected to an earlier chapter.
%
Giving every robot a local runtime instance, with localisation decentralised through radio ranging, reaches the monolithic end of the deployment space,
at which point a policy has two endpoints to choose between rather than one.
%
Injecting communication delays into the simulation engine gives the mixed deployments between those endpoints a cost,
which is the precondition for a policy of \cref{chap:deployment-and-reconfiguration} or \cref{chap:heterogeneous-gnn} to have anything to optimise.
%
Repeating the perturbation experiments under controlled conditions, rather than by hand in a crowded room, turns this chapter's qualitative observation into a measured convergence time.
%
And generating the collective behaviours from a natural-language specification, as done for ScaFi in \cref{chap:language-based-macroprogramming},
would put the whole chain from specification to physical actuation under test at once.
%
The four are taken up again in \cref{chap:conclusions}.
